\spread{Can reading a fact help a model write a sentence?}
\noindent A question about Chopin requires the writer to distinguish two dates on the same page. We trained a reader to return a strip of ink and asked a separate writer to turn that evidence into a complete answer.
\platehere{Figure_168_reading_before_writing}{A date can be read correctly and still disappear in the reply}{The two panels show the first exact sentence gained and lost when the writer received the predicted fact and its surrounding region, compared with the full-page writer. Each question, intermediate fact and model reply is original ink. The year in panel b was available to the writer when it omitted the date.}
\begin{multicols}{2}
\section*{Independent reading and writing}
The reader saw 39,057 English, simplified Chinese and traditional Chinese extractive questions. For a contiguous span from $i$ to $j$ within a source region its score was $u_i+v_j+b_{j-i+1}$. The two endpoints depend on the question and page images. A softmax over legal spans supplies the probability of an annotated fact by summing its accepted occurrences. Enumeration on three real pages agreed with all 31 parameter gradients to within $5.4\times10^{-7}$ in relative norm. Both the training objective and the final choice operate on visible glyphs.
\section*{What reaches the writer}
Three 865,605-parameter writers learned from the same initialization and the same 101,364 human answer wordings over four passes. One read the full page. A second learned from the annotated fact alone and received the reader's prediction at test. The third saw the predicted strip and its surrounding region in training and at test. The 116,547-parameter reader was trained on disjoint documents and queries. Independent checks found zero shared question images or ordered pages across fitting partitions and the new confirmation set. Content selection has an antecedent in neural summarization (\href{https://aclanthology.org/D18-1443/}{Gehrmann et al.}).
\section*{Completing the claim in English}
On 192 fresh questions, exact stopped answer images numbered four with the full page, ten with the fact alone and nine with the fact and region. Only the third writer answered both questions on any of the 96 passages, once. With a human fact the fact-only model wrote 41 exact answers. When the reader independently returned the exact annotated fact (22 cases), the writer completed ten reference statements. The other twelve answers lost part of the proposition or expressed it in different words. None of the remaining 170 cases yielded an exact statement.
\end{multicols}
\newpage
\spread{What survives a change of question?}
\noindent Two questions about the British Regency ask for an artistic period and a date range. Both models saw the same page each time.
\platehere{Figure_169_following_a_pictured_fact}{The reader selects the same Chinese phrase for two different requests}{The first question asks which artistic movement characterized the period. The reader and conditioned writer both return the neoclassical period. The following question asks for the years of the Regency. The reader repeats the artistic movement, and the conditioned writer reproduces it. The first listed human reference gives the full date range.}
\begin{multicols}{2}
\section*{Changing the requested relation}
On 192 new Chinese reading questions, exact answers rose from two for direct writing to 29 for the fact-only model and 13 for the fact-plus-region model. The respective numbers of both-correct question pairs were zero, six and one out of 96. The span reader alone also gave 29 accepted Chinese answers and six correct pairs. In this extractive task the best system preserved the reader's answer. With the human answer supplied, the fact-only model reproduced 185 of 192 reference images. When given a sibling question's predicted fact on the same page, it reproduced the original answer only three times.
\section*{Which changes become invisible?}
For a fixed question a deterministic writer must produce the same reply whenever its entire evidence input is identical. If the accepted answers for the two pages are disjoint, success on both is then impossible. A checked Lean theorem bounds the number of correct pairs by the number free of such collisions. Of 64 English and 190 Chinese substitutions, 38 and 92 produced identical evidence for the fact-only writer. All had disjoint reference images. Passing the region alongside the fact lowered the counts to 27 and 30. Complete edit-pair successes nevertheless fell from four to two in English and nine to zero in Chinese. Distinguishability is necessary but insufficient.
\section*{Where the reader loses the fact}
With the full-page scores fixed, restricting search to a marked region recovered 44 English and 37 Chinese facts. Giving a correct start instead recovered 128 and 92. Across all pages, 122 English and 119 Chinese accepted facts lay within the highest 32 span scores. The ordinary first choices recovered 22 and 29. All 384 CPU location choices reproduced those made on GPU. We then ran the fixed writer on every one of the 12,288 candidate positions. Even choosing with knowledge of the reference could yield at most 35 exact English and 115 exact Chinese replies. In English, 25 questions had an attainable answer that the reader missed and 157 had none. In Chinese those counts were 86 and 77. All 384 first choices reproduced the saved answer pixels and stopping decisions.
\end{multicols}
\newpage
\spread{What improves as the model learns to write?}
\noindent Six saved states allow the sequence of changes in the writer to be read against a fixed visual reader. Counting complete answer images distinguishes learning a fact from learning how to express it.
\platehere{Figure_170_learning_to_state_a_fact}{Learning curves for independent, freely generated answers}{Each point counts stopped replies on the same 64 development questions per language. The reader was fixed. The upper panels require the complete reference image. The lower panels count the presence of a requested fact. The 384 confirmation questions were reserved separately.}
\begin{multicols}{2}
\section*{From a fact to its relation}
After 5,000 updates the fact-only writer gave 18 exact Chinese answers on development. The count remained 18 at the last three measurements. In English it rose from one at 5,000 updates to two at 7,500 and stayed there. On confirmation, the paired-passage bootstrap 95\% intervals for the gain over direct writing were 1.0--5.7 percentage points in English and 8.9--19.8 in Chinese. Neither intervention improved correct natural question pairs and edited pairs in both languages as required. The test includes one initialization per arm. A blinded automatic judge accepted only ten of sixteen human references on edited English pages and five of sixteen on edited Chinese pages. Its failed calibration leaves semantic accuracy beyond exact matches unresolved.
\section*{Time, accuracy and transfer}
The reader and fact-only writer together have 982,152 parameters. Fitting took 52 and 171 seconds, respectively, on one RTX 4090 D. On 16 fresh questions per language, this pipeline and a local Q4 Qwen3-VL-8B used the same pictured content. The two systems returned one and four normalized exact English statements, and five each in Chinese. The English fact occurred in three ILM and thirteen Qwen answers. Generation used 275 and 9,713 J of device energy over all 32 questions. The difference in answer quality prevents a like-for-like efficiency estimate.
\section*{Limits of the learned skill}
Across 152 book prompts, none of the writers produced the reference eight-glyph prefix. The fact-only pipeline returned the requested historical specimen on one of 130 Chinese prompts and none of 130 English prompts. On a separate Chinese catalogue test the direct model produced 30 exact images in 96 trials and the fact-only system six. These losses restrict the evidence for transfer from reading short facts to longer answers and historical writing.
\end{multicols}
\newpage
